\section{Prove the raster extent before publishing pixels}
Once width \(w\) and height \(h\) satisfy the bounds, the pixel count is \(p=wh\).
This exercise has at most 4096 pixels. Each PPM pixel needs three bytes, giving
\(r=3p\); each framebuffer pixel occupies four bytes, giving \(f=4p\).
At the largest dimensions these are 12288 and 16384 bytes. The file arena needs room
for the separately bounded header as well as the raster.

Let \(o\) be the saved raster offset and \(n\) the actual file byte count. Establish
\(o\leq n\), then compare \(n-o\) with \(r\). Subtracting only after proving the order
avoids turning a short input into a huge unsigned remainder. Equality proves both that
the last pixel is present and that this single-image input has no trailer.

Validate the header, extent and destination capacity before the first framebuffer
store. Otherwise a late error leaves an apparently updated display containing a mixture
of old and new rows. The input byte buffer can be partial while a file read is in
progress; the visible framebuffer should change only after the file-level checks have
succeeded. These are separate publication boundaries, with different promises.

For the two-by-two fixture, calculate \(p=4\), \(r=12\), \(f=16\), \(o=11\) and
\(n=23\). Now shorten the file by one byte. The header remains plausible, but
\(n-o=11\) cannot satisfy the required raster extent. No pixel store should occur.
